%% ma: L12-B14-0013
%% lop: 12
%% chuong: 5
%% bai: 14
%% dang: 12.14.12
%% loai: TN4
%% muc: VD
%% dap_an: "A"
%% nguon: "(NBV)-ĐỀ SỐ 19-MỖI NGÀY 1 ĐỀ THI 2025.pdf (D19, MỖI NGÀY 1 ĐỀ - Đề số 19), Phần I Câu 12"
%% kien_thuc: "Bài 14 (phương trình mặt phẳng, tích có hướng của hai vectơ, mặt phẳng vuông góc với mặt phẳng cho trước)"
%% trang_thai: chinh-thuc
%% ghi_chu: "Đáp án gốc A đúng (giải độc lập: n_Q = [AB, n_P] = (0;8;12), (Q): 2y + 3z - 11 = 0, a + b + c = 5). Đã duyệt 19/10/2026."
\begin{ex}
Trong không gian với hệ trục tọa độ $Oxyz$, cho hai điểm $A(2;4;1)$, $B(-1;1;3)$ và mặt phẳng $(P)\colon x-3y+2z-5=0$. Mặt phẳng $(Q)$ đi qua hai điểm $A$, $B$ và vuông góc với mặt phẳng $(P)$ có phương trình dạng $ax+by+cz-11=0$. Khẳng định nào sau đây là đúng?
\choice{\True $a+b+c=5$}{$a+b+c=15$}{$a+b+c=-5$}{$a+b+c=-15$}
\loigiai{Ta có $\overrightarrow{AB}=(-3;-3;2)$ và vectơ pháp tuyến của $(P)$ là $\vec n_P=(1;-3;2)$. Vì $(Q)\perp(P)$ nên $\vec n_P$ là một vectơ chỉ phương của $(Q)$; ngoài ra $\overrightarrow{AB}$ cũng là một vectơ chỉ phương của $(Q)$. Suy ra $(Q)$ có vectơ pháp tuyến $\vec n_Q=[\overrightarrow{AB},\vec n_P]=(0;8;12)=4(0;2;3)$.

Phương trình $(Q)$ đi qua $A(2;4;1)$: $2(y-4)+3(z-1)=0\Leftrightarrow2y+3z-11=0$. Vậy $a=0$, $b=2$, $c=3$ và $a+b+c=5$. Chọn A.}
\end{ex}
